\chapter{Desarrollo}\label{Chap2}

\section{Marco teórico}

[Texto de ejemplo. Presente aquí los conceptos y definiciones necesarios
para entender el trabajo. Cada ecuación debe ir seguida de la definición
de sus símbolos, como en el ejemplo siguiente.]

\begin{equation}
y = ax^2 + bx + c
\label{eq:ejemplo}
\end{equation}

donde $y$ es la variable dependiente, $x$ es la variable independiente, y
$a$, $b$ y $c$ son coeficientes constantes. La ecuación~\eqref{eq:ejemplo}
se cita en el texto con su número, así: la ecuación~\eqref{eq:ejemplo}
describe una parábola.

\section{Metodología}

[Texto de ejemplo. Describa aquí cómo se realizó el trabajo: métodos,
materiales y procedimiento. Los conceptos de \cite{libro-ejemplo} sirven
de referencia para esta sección.]

\begin{figure}[htbp]
	\centering
	\begin{tikzpicture}[scale=0.9]
		\draw[->] (-2.5,0) -- (2.5,0) node[right] {$x$};
		\draw[->] (0,-0.5) -- (0,4.5) node[above] {$y$};
		\draw[domain=-2:2,smooth,variable=\x,thick] plot ({\x},{\x*\x});
	\end{tikzpicture}
	\caption[Figura de ejemplo]{Figura de ejemplo: parábola $y = x^2$.}
	\label{fig:ejemplo}
\end{figure}

La figura~\ref{fig:ejemplo} muestra un ejemplo de figura con su pie y su
etiqueta: toda figura debe introducirse y comentarse en el texto.

\section{Resultados de ejemplo}

[Texto de ejemplo. Presente aquí los resultados principales y compárelos
con lo reportado en la literatura, por ejemplo en
\cite{articulo-ejemplo}.]

\begin{table}[htbp]
	\centering
	\caption[Tabla de ejemplo]{Tabla de ejemplo: datos ilustrativos.}
	\label{tab:ejemplo}
	\begin{tabular}{lcc}
	\hline
	Caso & Valor A & Valor B \\
	\hline
	1 & 1.0 & 2.0 \\
	2 & 1.5 & 2.5 \\
	3 & 2.0 & 3.0 \\
	\hline
	\end{tabular}
\end{table}

La tabla~\ref{tab:ejemplo} resume los datos de ejemplo: toda tabla debe
introducirse y comentarse en el texto, igual que las figuras.
